% ============================================================
% Chapter 1: Introduction
% ============================================================
\chapter{Introduction}
\label{ch:introduction}

\section{Research Motivation}
\label{sec:motivation}

Optimisation lies at the heart of countless scientific and engineering
endeavours, from tuning hyperparameters of deep neural networks to designing
novel materials with desired properties~\cite{bibid}.  In many of these
applications, the objective function is expensive to evaluate, noisy, and
lacks gradient information---characteristics that render classical gradient-based
optimisation methods inapplicable.  Bayesian optimisation
(\gls{bo}) has emerged as the method of choice for such black-box
optimisation problems, offering a principled framework that balances exploration
and exploitation through surrogate modelling and acquisition function
maximisation~\cite{einstein1905}.

Despite its success in low-dimensional settings ($d \leq 20$), the
effectiveness of standard BO degrades sharply in higher dimensions due to the
curse of dimensionality~\cite{goossensLaTeXCompanion1993}.  The volume of the
search space grows exponentially with dimension, rendering the surrogate model's
predictions unreliable across most of the domain.  This limitation represents a
critical bottleneck, as many real-world problems---such as molecular property
optimisation in computational chemistry or aerodynamic shape design using
\gls{cfd} simulations---inherently involve tens to hundreds of
parameters.

\section{Research Objectives}
\label{sec:objectives}

The primary goal of this thesis is to develop a scalable \gls{bo} framework
that maintains sample efficiency in dimensions ranging from 50 to 500.  The
specific objectives are:

\begin{enumerate}
    \item \textbf{Objective 1:} Design an adaptive subspace identification
    algorithm that learns the intrinsic low-dimensional structure of the
    objective function without imposing rigid manifold assumptions.

    \item \textbf{Objective 2:} Derive information-theoretic acquisition
    functions that naturally adapt to varying local dimensionality across the
    search space.

    \item \textbf{Objective 3:} Establish theoretical convergence guarantees
    for the proposed methods under mild regularity conditions.

    \item \textbf{Objective 4:} Validate the framework on standard benchmark
    functions and demonstrate its practical utility on real-world engineering
    problems.
\end{enumerate}

\section{Contributions}
\label{sec:contributions}

This thesis makes the following contributions to the field of high-dimensional
optimisation:

\begin{enumerate}
    \item A novel adaptive subspace BO (\gls{asbo}) algorithm with provable
    convergence guarantees in dimensions up to 500, significantly extending the
    practical range of BO methods.

    \item A family of information-theoretic acquisition functions---including
    predictive entropy search (PES) and its variants---that adapt to varying
    local dimensionality, providing stronger exploration guarantees than
    \gls{ucb}-based approaches.

    \item A comprehensive theoretical analysis of convergence rates, including
    bounds on simple regret that depend on the intrinsic dimensionality rather
    than the ambient dimension.

    \item An open-source Python library (\texttt{highdimbo}) implementing the
    proposed methods, with seamless integration into the
    \href{https://github.com/pytorch/botorch}{BoTorch} framework.

    \item Extensive empirical evaluation on 12 synthetic benchmarks and 3
    real-world engineering problems, demonstrating state-of-the-art performance.
\end{enumerate}

\section{Thesis Outline}
\label{sec:outline}

The remainder of this thesis is organised as follows:

\begin{itemize}
    \item \Cref{ch:literature-review} provides a systematic review of existing
    approaches to high-dimensional BO, including dimensionality reduction
    techniques, trust-region methods, and information-theoretic approaches.

    \item \Cref{ch:methodology} presents the formal problem definition and
    details the proposed ASBO algorithm, including the adaptive subspace
    identification module, local surrogate modelling, and the family of
    information-theoretic acquisition functions.

    \item \Cref{ch:experiments} describes the experimental setup and presents
    results on synthetic benchmark functions and real-world engineering problems.

    \item \Cref{ch:analysis} provides a detailed statistical analysis of the
    experimental results, including ablation studies, scalability analysis, and
    failure mode analysis.

    \item \Cref{ch:conclusion} summarises the key findings, discusses
    limitations, and outlines directions for future research.
\end{itemize}

\section{Thesis Publications}
\label{sec:publications}

The results presented in this thesis are based on the following publications:

\begin{enumerate}
    \item \textbf{J. Smith, A. Turing, and D. Knuth.}
    ``Adaptive Subspace Bayesian Optimisation for High-Dimensional Black-Box
    Functions.''
    \textit{Advances in Neural Information Processing Systems (NeurIPS)}, 2025.

    \item \textbf{J. Smith and A. Turing.}
    ``Information-Theoretic Acquisition Functions for Scalable Bayesian
    Optimisation.''
    \textit{International Conference on Machine Learning (ICML)}, 2025.

    \item \textbf{J. Smith, D. Knuth, and G. Hopper.}
    ``\texttt{highdimbo}: An Open-Source Library for High-Dimensional Bayesian
    Optimisation.''
    \textit{Journal of Machine Learning Research (JMLR)}, 2026.
\end{enumerate}

\section{Notation and Conventions}
\label{sec:notation}

Throughout this thesis, we adopt the following notation conventions.  Vectors
are denoted by bold lowercase letters ($\bx, \by, \bw$) and matrices by bold
uppercase letters ($\bPhi, \bSigma$).  The $i$-th element of a vector $\bx$ is
denoted $x_i$, and the $(i,j)$-th element of a matrix $\bA$ is $A_{ij}$.  We
use $\R^d$ to denote the $d$-dimensional real vector space and $\cN(\bmu, \bSigma)$
to denote a multivariate Gaussian distribution with mean $\bmu$ and covariance
$\bSigma$.  The notation $\cO(\cdot)$ denotes big-O complexity.  Probabilistic
notation follows standard conventions: $\Pr(\cdot)$ denotes probability,
$\E[\cdot]$ denotes expectation, and $\Var(\cdot)$ denotes variance.  Key
notation is summarised in \cref{tab:notation}.

\begin{table}[htbp]
    \centering
    \caption{Summary of mathematical notation used in this thesis.}
    \label{tab:notation}
    \begin{tabular}{@{}ll@{}}
        \toprule
        \textbf{Symbol} & \textbf{Description} \\
        \midrule
        $d$ & Ambient dimensionality \\
        $k$ & Intrinsic dimensionality (subspace) \\
        $n$ & Number of function evaluations \\
        $f(\bx)$ & Objective function \\
        $\cO$ & Observation space \\
        $\cX$ & Input domain, $\cX \subseteq \R^d$ \\
        $\mu_n(\bx)$ & Posterior mean at iteration $n$ \\
        $\sigma_n(\bx)$ & Posterior std.\ at iteration $n$ \\
        $\alpha_n(\bx)$ & Acquisition function at iteration $n$ \\
        $\bZ \in \R^{k \times d}$ & Subspace projection matrix \\
        $\cN(\bmu, \bSigma)$ & Gaussian distribution \\
        $\KL(\cdot \| \cdot)$ & Kullback--Leibler divergence \\
        \bottomrule
    \end{tabular}
\end{table}
